\subsection{APPROXIMATION OF TOPOLOGICAL GROUPS}

Let G be a group and let \((\mathrm{H}_{\alpha})_{\alpha\in\mathbf{I}}\) be a family of normal subgroups of G such that \(H_{\alpha}\supset H_{\beta}\) whenever \(\alpha\leqslant\beta\) . For each \(\alpha\in I\) let \(G_{\alpha}=G/H_{\alpha}\) , and for \(\alpha\leqslant\beta\) let \(f_{\alpha\beta}\) be the canonical homomorphism \(G/H_{\beta}\to G/H_{\alpha}\) , which therefore maps a coset T of \(H_{\beta}\) in G to the coset \(TH_{\alpha}\) of \(H_{\alpha}\) in G. Clearly \((G_{\alpha},f_{\alpha\beta})\) is an inverse system of groups; the elements of \(\tilde{G}=\varprojlim G_{\alpha}\) are families \((T_{\alpha})_{\alpha\in\mathbf{I}}\) , where \(T_{\alpha}\) is a coset of \(H_{\alpha}\) in G for each \(\alpha\) , and \(T_{\alpha}\supset T_{\beta}\) whenever \(\alpha\leqslant\beta\) . The mapping \(i:s\to(sH_{\alpha})\) of G into \(\tilde{G}\) is the inverse limit of the canonical homomorphisms \(G\to G/H_{\alpha}\) and is therefore a homomorphism of G into \(\tilde{G}\) , and the inverse image under i of an element \((T_{\alpha})\in\tilde{G}\) is equal to \(\bigcap_{\alpha\in I}T_{\alpha}\) . The kernel of i is therefore \(\bigcap_{\alpha\in I}H_{\alpha}\) , and the image of i consists of all families \((T_{\alpha})\in\tilde{G}\) whose intersection is non-empty.

Now suppose that G is a topological group; if we give each \(G_{\alpha} = G/H_{\alpha}\) the quotient topology, it is clear that \((\mathbf{G}_{\alpha}, f_{\alpha\beta})\) is an inverse system of topological groups, and that \(i : G \to \tilde{G}\) is a continuous homomorphism.

PROPOSITION 2. Let G be a topological group and let \((\mathrm{H}_{\alpha})_{\alpha\in I}\) be a family of normal subgroups of G such that \(H_{\alpha}\supset H_{\beta}\) whenever \(\alpha\leqslant\beta\) , and which satisfy the following condition:

(AP) For each \(\alpha \in I\) , \(H_{\alpha}\) is closed in G and every neighbourhood of the identity element \(e\) in G contains one of the \(H_{\alpha}\) (in other words, the filter base formed by the \(H_{\alpha}\) converges to \(e\) ).

Then the mapping \(i: \mathbf{G} \to \tilde{\mathbf{G}} = \varprojlim \mathbf{G} / \mathbf{H}_{\alpha}\) is a strict morphism of \(\mathbf{G}\) onto \(i(\mathbf{G})\) ; \(\tilde{\mathbf{G}}\) is Hausdorff and \(i(\mathbf{G})\) is dense in \(\tilde{\mathbf{G}}\) ; and finally the kernel of \(i\) is the closure of \(\{e\}\) in \(\mathbf{G}\) . If in addition one of the \(\mathbf{H}_{\alpha}\) is complete, then \(i\) is surjective.

Clearly the \(G_{\alpha} = G/H_{\alpha}\) are Hausdorff (§ 2, no. 5, Proposition 13), hence so is \(\tilde{G}\) (since it is a subspace of \(\prod_{\alpha \in I} G_{\alpha}\) ). The kernel H of i is the intersection of the \(H_{\alpha}\) and is therefore a closed subgroup of G. Since each neighbourhood of \(e\) contains some \(H_{\alpha}\) , it contains \(H\) and therefore {[}§ 3, no. 1, formula (1){]} \(H\) is the closure of \(\{e\}\) . Let us next show that \(i(G)\) is dense in \(\tilde{G}\) . Let \(f_{\alpha}\) be the canonical mapping \(\tilde{G} \to G_{\alpha}\) , which is the restriction to \(\tilde{G}\) of the projection \(pr_{\alpha}\) ; \(\varphi_{\alpha} = f_{\alpha} \circ i\) is the canonical mapping \(G \to G / H_{\alpha}\) . If \(U\) is any non-empty open set of \(\tilde{G}\) , then there is an index \(\alpha \in I\) and a non-empty open set \(U_{\alpha}\) in \(G_{\alpha}\) such that \(\overline{f}_{\alpha}^{-1}(U_{\alpha}) \subset U\) (Chapter I, § 4, no. 4, Proposition 9); therefore

\[
\overline {{i}} ^ {-1} (\mathbf {U}) \supset \overline {{\varphi}} _ {\alpha} ^ {-1} (\mathbf {U} _ {\alpha});
\]

but since \(\varphi_{\alpha}\) is surjective, \(\overline{i}^{-1}(\mathbf{U})\) is not empty, and therefore \(i(\mathbf{G}) \cap \mathbf{U} \neq \emptyset\) . To see that \(i\) is a strict morphism of \(\mathbf{G}\) onto \(i(\mathbf{G})\) , consider a neighbourhood \(\mathbf{V}\) of \(e\) in \(\mathbf{G}\) . There is a neighbourhood \(\mathbf{W}\) of \(e\) in \(\mathbf{G}\) such that \(\mathbf{W}^2 \subset \mathbf{V}\) , and an index \(\alpha \in \mathbf{I}\) such that \(\mathbf{H}_{\alpha} \subset \mathbf{W}\) ; it follows that \(\mathbf{V}\) contains \(\mathrm{WH}_{\alpha} = \overline{\varphi}_{\alpha}^{-1}(\varphi_{\alpha}(\mathbf{W})) = \overline{i}^{-1}(\overline{f}_{\alpha}^{-1}(\varphi_{\alpha}(\mathbf{W})))\) . Since \(\overline{f}_{\alpha}^{-1}(\varphi_{\alpha}(\mathbf{W}))\) is a neighbourhood of the identity element in \(\tilde{\mathbf{G}}\) , the result follows.

Finally, suppose that there is an index \(\gamma \in I\) such that \(H_{\gamma}\) is complete. To show that \(i\) is surjective it is enough to prove that every family \((T_{\alpha})_{\alpha \in I} \in \tilde{G}\) has a non-empty intersection. Since \(T_{\gamma}\) is obtained by translation from \(H_{\gamma}\) , it is a complete subspace of \(G\) (with respect to both the right and the left uniformities). Moreover, since every neighbourhood \(U\) of \(e\) in \(G\) contains one of the \(H_{\alpha}\) , the corresponding set \(T_{\alpha}\) is \(U_{d-}\) (or \(U_{s-}\) ) small, and hence the set of \(T_{\alpha}\) contained in \(T_{\gamma}\) is a Cauchy filter base; therefore it converges in \(T_{\gamma}\) , and since the \(T_{\alpha}\) are closed in \(G\) (since they are translations of the \(H_{\alpha}\) ), their intersection is not empty.

Q.E.D.

COROLLARY 1. If the condition (AP) is satisfied and if in addition the (Hausdorff) groups \(\mathbf{G} / \mathbf{H}_{\alpha}\) are complete, then the group \(\mathbf{G}\) has a Hausdorff completion which can be identified with \(\tilde{\mathbf{G}}\) ; and the mapping \(i: \mathbf{G} \to \tilde{\mathbf{G}}\) is then identified with the canonical mapping (§ 3, no. 3, Proposition 5).

For \(\tilde{\mathbf{G}}\) is then complete (no. 2), and Proposition 2 shows that \(i(\mathbf{G})\) is isomorphic with the Hausdorff group associated with \(\mathbf{G}\) ; since it is dense in \(\tilde{\mathbf{G}}\) , the corollary follows ( \(\S 3\) , no. 3, Proposition 5). In particular:

COROLLARY 2. Let G be a group and let \((\mathrm{H}_{\alpha})\) be a family of normal subgroups of G, directed with respect to the relation \(\mathrm{H}_{\alpha} \supset \mathrm{H}_{\beta}\) . If we endow G with the group topology for which the \(\mathrm{H}_{\alpha}\) form a fundamental system of neighbourhoods of the identity element \(e\) , then the Hausdorff group associated with G is isomorphic to \(\mathrm{G}_{1} = \mathrm{G} / \left( \bigcap_{\alpha} \mathrm{H}_{\alpha} \right)\) ; \(\mathrm{G}_{1}\) has a completion \(\hat{\mathrm{G}}_{1} = \tilde{\mathrm{G}}\) ; and the canonical mapping \(\mathrm{G}_{1} \to \hat{\mathrm{G}} = \varprojlim \mathrm{G} / \mathrm{H}_{\alpha}\) extends to an isomorphism of \(\hat{\mathrm{G}} = \hat{\mathrm{G}}_{1}\) onto \(\tilde{\mathrm{G}}\) .

The subgroup \(H_{\alpha}\) of G is open and therefore also closed ( \(\S\) 2, no. 1, Corollary to Proposition 4), and \(G/H_{\alpha}\) is discrete ( \(\S\) 2, no. 6, Proposition 18); hence the conditions of Corollary 1 are satisfied.

For the remainder of this sub-section we shall assume that G is Hausdorff and that \((\mathbf{H}_{\alpha})\) is a family of compact normal subgroups of G, which is directed with respect to the relation \(H_{\alpha} \supset H_{\beta}\) and which satisfies the condition (AP); by virtue of Proposition 2, the mapping

\[
i: \mathrm{G} \rightarrow \tilde {\mathrm{G}} = \varprojlim \mathrm{G} / \mathrm{H} _ {\alpha}
\]

is then an isomorphism of topological groups which permits us to identify G and \(\tilde{G}\) . We denote by \(f_{\alpha}\) the canonical mapping \(G \rightarrow G/H_{\alpha}\) .

LEMMA 1. Under the hypotheses of Proposition 2, if E is any closed subset of G we have \(\mathbf{E} = \bigcap_{\alpha} \mathrm{EH}_{\alpha}\) .

For E is the intersection of the sets EV where V runs through the neighbourhood filter of e {[}§ 3, no. 1, formula (1){]}, and every neighbourhood of e contains an \(H_{\alpha}\) ; whence the result, since \(E \subset EH_{\alpha}\) .

PROPOSITION 3. Suppose that G is Hausdorff and that the \(H_{\alpha}\) are compact and satisfy (AP).

\begin{enumerate}
\def\labelenumi{\alph{enumi})}
\item
  Let L be a closed subgroup of G; then, for each \(\alpha\in I\) , the subgroup \(L_{\alpha}=f_{\alpha}(L)\) of \(G_{\alpha}=G/H_{\alpha}\) is closed, and the isomorphism i of G onto \(\varprojlim G_{\alpha}\) gives by restriction an isomorphism of L onto \(\varprojlim L_{\alpha}\) . If also L is normal in G, then \(L_{\alpha}\) is normal in \(G_{\alpha}\) for each \(\alpha\in I\) , and by passing to the quotients, i induces an isomorphism of G/L onto \(\varprojlim G_{\alpha}/L_{\alpha}\) .
\item
  Conversely, for each \(\alpha \in I\) let \(L_{\alpha}\) be a closed subgroup of \(G_{\alpha}\) , such that \(L_{\alpha} = f_{\alpha \beta}(L_{\beta})\) whenever \(\alpha \leqslant \beta\) . Then there is a unique closed subgroup \(L\) of \(G\) such that \(L_{\alpha} = f_{\alpha}(L)\) for each \(\alpha \in I\) ; and if in addition \(L_{\alpha}\) is normal in \(G_{\alpha}\) for each \(\alpha \in I\) , then \(L\) is normal in \(G\) .
\item
  Since \(\mathbf{H}_{\alpha}\) is compact, \(\mathrm{LH}_{\alpha}\) is closed in \(\mathbf{G}\) ( \(\S 4\) , no. 1, Proposition 1, Corollary 1) and therefore \(\mathbf{L}_{\alpha}\) is closed in \(\mathbf{G}_{\alpha}\) . Since \(i\) identifies the topological groups \(\mathbf{G}\) and \(\varprojlim \mathbf{G}_{\alpha}\) , and since \(\varprojlim \mathbf{L}_{\alpha}\) may be identified with a (topological) subgroup of \(\varprojlim \mathbf{G}_{\alpha}\) , \(i\) identifies the subgroup \(\bigcap_{\alpha} \mathrm{LH}_{\alpha}\) of \(\mathbf{G}\) with \(\varprojlim \mathbf{L}_{\alpha}\) , and to prove the first assertion it is enough to remark that \(\mathbf{L} = \bigcap_{\alpha} \mathrm{LH}_{\alpha}\) by Lemma 1. On the other hand, if \(\mathbf{L}\) is normal, then for each \(\alpha \in \mathbf{I}\) the mapping \(f_{\alpha}': \mathbf{G} / \mathbf{L} \to \mathbf{G}_{\alpha} / \mathbf{L}_{\alpha}\) induced by \(f_{\alpha}\) is a surjective strict morphism ( \(\S 2\) , no. 8, Remark 3), whose kernel is the compact normal subgroup \(\mathrm{H}_{\alpha} \mathrm{L} / \mathrm{L}\) of \(\mathrm{G} / \mathrm{L}\) , the canonical image of the compact subgroup \(\mathrm{H}_{\alpha}\) of \(\mathbf{G}\) . Since the subgroups \(\mathrm{H}_{\alpha} \mathrm{L} / \mathrm{L}\) of \(\mathrm{G} / \mathrm{L}\) satisfy condition (AP) (§ 2, no. 8, Proposition 24), and since G/L is Hausdorff, the last assertion of a) is a consequence of Proposition 2.
\item
  Let \(f_{\alpha \beta}'\) be the restriction of \(f_{\alpha \beta}\) to \(\mathbf{L}_{\beta} (\alpha \leqslant \beta)\) . Then \((\mathbf{L}_{\alpha}, f_{\alpha \beta}')\) is an inverse system of topological groups, whose inverse limit \(\mathbf{L}\) can be identified with the subgroup \(\mathbf{G} \cap \prod_{\alpha} \mathbf{L}_{\alpha}\) of \(\mathbf{G}\) . By hypothesis, \(f_{\alpha \beta}'\) is surjective and its kernel is the compact subgroup \(f_{\beta}(\mathrm{H}_{\alpha}) \cap \mathrm{L}_{\beta}\) of \(\mathbf{L}_{\beta}\) ; consequently (Chapter I, § 9, no. 6, Corollary 1 to Proposition 8) we have \(\mathbf{L}_{\alpha} = f_{\alpha}(\mathbf{L})\) for all \(\alpha \in \mathbf{I}\) . If \(\mathbf{L}'\) is another closed subgroup of \(\mathbf{G}\) such that \(f_{\alpha}(\mathbf{L}') = \mathbf{L}_{\alpha}\) for all \(\alpha \in \mathbf{I}\) , then \(\mathbf{L}'\mathbf{H}_{\alpha} = \overline{f}_{\alpha}^{-1}(\mathbf{L}_{\alpha})\) , whence (Lemma 1) \(\mathbf{L}' = \bigcap_{\alpha} \mathbf{L}'\mathbf{H}_{\alpha} = \bigcap_{\alpha} \overline{f}_{\alpha}^{-1}(\mathbf{L}_{\alpha}) = \mathbf{L}\) . Finally, the last assertion of b) follows from the formula \(\mathbf{L} = \bigcap_{\alpha} \overline{f}_{\alpha}^{-1}(\mathbf{L}_{\alpha})\) , since the \(\overline{f}_{\alpha}^{-1}(\mathbf{L}_{\alpha})\) are now normal subgroups of \(\mathbf{G}\) .
\end{enumerate}

PROPOSITION 4. Suppose that G is Hausdorff, the \(H_{\alpha}\) are compact and that (AP) is satisfied. If \(C_{\alpha}\) is the identity component of \(G_{\alpha} = G/H_{\alpha}\) , then the identity component C of G can be identified with \(\varprojlim C_{\alpha}\) , and we have \(f_{\alpha}(C) = C_{\alpha}\) . The proposition is a consequence of the following lemma:

LEMMA 2. Let G be a Hausdorff topological group and let H be a compact normal subgroup of G, and \(\varphi\) the canonical mapping \(\mathbf{G} \to \mathbf{G} / \mathbf{H}\) . If C is the identity component of G, then \(\varphi(\mathbf{C})\) is the identity component of \(\mathbf{G}' = \mathbf{G} / \mathbf{H}\) .

Once this lemma is established, we shall have \(f_{\alpha}(\mathbf{C}) = \mathbf{C}_{\alpha}\) for each \(\alpha \in \mathbf{I}\) , and since \(\mathbf{C}\) is a closed subgroup of \(\mathbf{G}\) (§ 2, no. 2, Proposition 7) it suffices to apply Proposition 3a).

To prove the lemma, notice first that if \(C'\) is the component of the identity element \(e'\) in \(G'\) , then \(\varphi(C) \subset C'\) since \(\varphi(C)\) is connected. Suppose that \(\varphi(C) \neq C'\) . Since C is a closed normal subgroup of G (§2, no. 2, Proposition 7), \(\varphi(C)\) is a normal subgroup of \(G'\) ; if \(\psi\) is the canonical mapping \(G' \to G'/\varphi(C)\) , then \(\psi(C')\) is connected and does not consist of the identity element alone; hence the identity component of \(G'/\varphi(C)\) does not consist of the identity element alone. But \(G'/\varphi(C)\) is isomorphic to \((G/H)/(HC/H)\) , hence to G/HC, and consequently also to \((G/C)/(HC/C)\) (§2, no. 7, Corollary to Proposition 22, and Proposition 20). Now, G/C is Hausdorff and totally disconnected (Chapter I, §11, no. 5, Proposition 9), and HC/C, being the canonical image of the compact normal subgroup H of G, is a compact subgroup of G/C. It is therefore sufficient to prove Lemma 2 under the additional hypothesis that G is totally disconnected, i.e.~\(C = \{e\}\) .

Suppose then that \(C' \neq \{e'\}\) ; replacing G by its subgroup \(\varphi^{-1}(C')\) , which is totally disconnected and contains H, we may suppose that \(G'\) is connected and does not consist of a single point.

Let \(\mathfrak{M}\) be the set of closed subgroups \(L\) of \(G\) such that \(LH = G\) . We shall show that the set \(\mathfrak{M}\) , ordered by the relation \(\supset\) , is inductive. Indeed, if \(\mathfrak{X}\) is a linearly ordered subset of \(\mathfrak{M}\) , then for each \(x \in G\) the set of sets \(xH \cap L\) with \(L \in \mathfrak{X}\) is a filter base formed of closed sets in the compact space \(xH\) ; hence the intersection of these sets is not empty, which shows that the intersection of the subgroups \(L \in \mathfrak{X}\) belongs to \(\mathfrak{M}\) . Applying Zorn's lemma, we see therefore that \(\mathfrak{M}\) has a minimal element \(L_0\) . Since \(H\) is compact, \(G/H = L_0H/H\) is isomorphic to \(L_0/(L_0 \cap H)\) ( \(\S 4\) , no. 1, Proposition 1, Corollary 3); since \(L_0\) is totally disconnected and \(L_0 \cap H\) is compact, we see that we may replace \(G\) by \(L_0\) ; in other words, we may suppose in addition that there is no closed subgroup \(L \neq G\) such that \(LH = G\) .

Now let F be the intersection of the neighbourhoods of the identity element which are both open and closed in G, and let us show that F is a closed subgroup of G. Clearly F is closed in G; hence it is enough to show that \(F^{-1}.F \subset F\) . But if \(x \in F\) and if V is an open and closed neighbourhood of e in G, then so is xV; for otherwise e would belong to the complement W of xV in G, which is again open and closed, and we should have \(x \notin W\) ; hence \(x \notin F\) , contrary to hypothesis. It follows that xF, the intersection of the sets xV as V runs through the open and closed neighbourhoods of e in G, contains F; in other words, \(x^{-1}F \subset F\) , which proves our assertion. Since G is totally disconnected and \(G \neq \{e\}\) , we have \(F \neq G\) . But if V is an open and closed neighbourhood of e in G, then VH is also open and closed in G (§ 4, no. 1, Corollary 1 to Proposition 1), hence \(\varphi(V)\) is both open and closed in G/H, and this implies that \(\varphi(V) = G/H\) by virtue of the hypothesis. We shall conclude from this that FH = G, which will give us the desired contradiction and complete the proof of the lemma. Indeed, for each \(x \in G\) , xH meets every neighbourhood V of e which is both open and closed, hence also meets the intersection F of these neighbourhoods, since the sets \(Vx \cap H\) form a filter base consisting of closed sets in the compact space xH.

Q.E.D.

Remark. If the subgroup \(H_{\alpha}\) is compact for some \(\alpha \in I\) , then \(H_{\beta}\) is compact for all \(\beta \geqslant \alpha\) , because it is a closed subgroup of \(H_{\alpha}\) . Since the set of all \(\beta \in I\) such that \(\beta \geqslant \alpha\) is cofinal in I, it makes essentially no difference, in the study of the group G, whether we suppose that one of the \(H_{\alpha}\) is compact or that all the \(H_{\alpha}\) are compact.

